\chapter{Die Parallelisierung}

\section{Zenon}

Fink liest die Paradoxien Zenons nicht als Sophismen, sondern als die erste systematische Dialektik von Raum, Zeit und Bewegung. Ihre Wurzel ist ein ungeklärtes \zitat{In-sein}. Dass ein Ding im Raum ist, versteht Zenon nach dem Muster des Gefäßes: Es besetzt einen Ort und keinen anderen. Dass ein Ding in der Zeit ist, versteht er nach demselben Muster: Es besetzt ein Jetzt und kein anderes. Fink nennt das ein \zitat{Doppelgeleise}. Auf beiden Gleisen läuft dieselbe Logik, und auf beiden entsteht dieselbe Schwierigkeit: Das Kontinuum ist vor seinen Teilen, aber das Muster des Gefäßes verlangt, dass es aus ihnen bestehe. Der fliegende Pfeil ruht, weil er in jedem Jetzt an einem Ort ist; Achilles erreicht die Schildkröte nicht, weil zwischen zwei Orten unendlich viele liegen.\footnote{Fink, Frühgeschichte, Kapitel zu Zenon.}

Die stillschweigende Voraussetzung, sagt Fink, ist, ein Ding könne nicht zugleich an zwei Orten und nicht am selben Ort zu zwei Zeiten sein. \zitat{In dieser massiven Formulierung ist das sicher richtig, aber Bewegung ist die seltsame Weise, wo ein Ding im selben Moment an zwei Orten ist, d.\,h. im Übergang zwischen zwei Orten sich befindet.}\footnote{Fink, Frühgeschichte, Kapitel zu Zenon.} Was hier auffällt, ist, dass der Übergang zeitlich beschrieben wird (\zitat{im selben Moment}) und räumlich gescheitert ist (\zitat{an zwei Orten}). Die Parallelisierung bricht genau da, wo sie sich bewähren sollte, in der Bewegung, und sie bricht auf der Seite des Raumes. Vom Raum her lässt sich der Übergang nicht denken, weil ein Ort ein Ort ist. Von der Zeit her lässt er sich denken, weil ein Jetzt nie nur ein Jetzt ist; was das heißt, wird sich an der Grenze zeigen, die das Jetzt ist. Das ist die erste Spur der Asymmetrie: Sie liegt nicht in einem Merkmal, das die Zeit hat und der Raum nicht, sondern darin, dass die Zeit den Übergang trägt, den der Raum zerlegt.

\section{Kant}

Reichenbach hat Kant vorgeworfen, Raum und Zeit \zitat{in einem gemeinsamen Kapitel} behandelt zu haben, so als wären es zwei Fälle desselben, die Anschauungsform.\footnote{Reichenbach, Philosophie der Raum-Zeit-Lehre, \S\,16.} Der Vorwurf trifft die transzendentale Ästhetik. Er trifft nicht die Analytik. Marko Martić hat gezeigt, dass Kant in der Analytik einen \emph{synthetischen} Vorrang der Zeit ansetzt: Die Schemata sind Zeitbestimmungen, die Reihe ist zeitlich, der Raum dagegen ist \zitat{ein Aggregat, aber keine Reihe}.\footnote{Kant, Kritik der reinen Vernunft, B\,439, nach Martić, Zeit und Raum.} Was aufgezählt werden soll, muss nacheinander aufgezählt werden, und das Nacheinander ist Zeit. Zugleich hält Kant an einer Gegenrichtung fest, die in der Widerlegung des Idealismus ausgesprochen wird: Die Bestimmung meiner Existenz in der Zeit setzt ein Beharrliches voraus, und das Beharrliche ist nicht in mir, sondern im Raum. In einer Reflexion heißt es, ohne Raum würde die Zeit selbst nicht als Größe vorgestellt werden. Martić folgt Reinhard Brandt in der Folgerung, \zitat{daß der Raum also die Zeit voraussetzt und nur als Zeitraum konzipierbar ist}, und hält zugleich die Umkehrung fest: Die Zeit wird nur am Raum zur Größe.\footnote{Martić, Zeit und Raum, Kapitel zum Schematismus und zur Widerlegung des Idealismus.} Das ist nicht mehr Parallelisierung. Es ist ein Wechselverhältnis, in dem die Zeit die Ordnung stiftet und der Raum das Maß ermöglicht. Es wird sich zeigen, dass Reichenbach dieselbe Struktur aus der Physik gewinnt.

\section{Minkowski}

Die vollständigste Gestalt der Parallelisierung ist die, die die Zeit zur weiteren Koordinate macht. Bei Minkowski wird sie als imaginäre Koordinate eingeführt, damit das Linienelement die Form einer euklidischen Länge annimmt. Ilse Schneider hat schon 1921 festgehalten, was diese Wendung bedeutet und was nicht: Raum und Zeit \zitat{verlieren die physikalische Gegenständlichkeit auch vereint, und nur durch die Materie kommt die physikalische Bedeutung hinein}.\footnote{Schneider, Das Raum-Zeit-Problem bei Kant und Einstein, Kapitel zur allgemeinen Relativitätstheorie.} Die Koordinaten sind Symbole, messbar ist das Linienelement. Damit ist die Zeit zwar der Form nach den Raumkoordinaten gleichgestellt, aber die Gleichstellung ist eine Gleichstellung von Symbolen, nicht von Sachen.

Reichenbach geht einen Schritt weiter. Für ihn ändert die Auffassung der Zeit als weitere Dimension nichts am Sondercharakter der Zeit. Das negative Vorzeichen, das in der metrischen Fundamentalform vor dem Zeitglied steht, ist nicht der Preis, den man für die Gleichstellung zahlt, sondern der mathematische Ausdruck dafür, dass sie nicht gelingt. Und er fügt hinzu, dass das Wesen der Zeit mit dieser Charakterisierung \zitat{noch nicht erschöpft} ist: Das Vorzeichen spricht nur eine Trennung von Raum und Zeit aus, nicht den eigentümlichen Charakter der Zeit.\footnote{Reichenbach, Raum-Zeit-Lehre, \S\,16 und \S\,30.} Ob man das Vorzeichen als solches schreibt oder es in eine imaginäre Einheit versteckt, ist Konvention. Dass eine der Koordinaten anders ist als die anderen, ist keine. An dieser Stelle zeigt sich am deutlichsten, was oben behauptet wurde: Die Parallelisierung legt die Asymmetrie an einen Ort, wo sie nicht auffällt. Sie schafft sie nicht ab.

\section{Hegel}

Hegel nennt Raum und Zeit \zitat{zwei Weisen des Außersichseins}: Der Raum ist das gleichgültige Außersichsein als ruhendes, die Zeit dasselbe in seiner Negativität, als Werden.\footnote{Hegel, Naturphilosophie, Nachschriften, nach Bonsiepen.} Das klingt nach Parallelität, und Bonsiepen hat gezeigt, dass Hegel den Raum tatsächlich in einer Weise denkt, die dem Sein des Parmenides näher steht als dem Werden: als Sichselbstgleichheit, als daseienden Äther. Aber die beiden Weisen sind nicht gleichrangig. Der Punkt, der im Raum die Negation des Raumes sein soll, hat im Raum keine Wirklichkeit; er wird sogleich zur Linie, die Linie zur Fläche. Erst in der Zeit hat die Negativität ihre Wirklichkeit. Deshalb der Satz aus der Nachschrift: \zitat{Die Zeit ist also erst der wahrhafte Punkt.} Und deshalb geht bei Hegel der Raum in die Zeit über, nicht die Zeit in den Raum. Die Parallelisierung ist bei ihm ein Anfang, der sich selbst aufhebt.

Das ist die Lage. Alle, die Zeit und Raum parallel behandelt haben, wissen um den Unterschied. Zenon scheitert an ihm, Kant baut ihn in die Analytik ein, Minkowski versteckt ihn im Vorzeichen, Hegel lässt ihn im Übergang hervortreten. Er wird nun an den drei Stellen herausgeholt, an denen er sich am wenigsten verbergen lässt: an der Grenze, an der Richtung, am Maß.
